% !TeX root = 01_intro_ml_polytech.tex
% !TeX encoding = UTF-8
% !TeX spellcheck = en_GB

\begin{frame}{The Jupyter Interface}{\includegraphics[height=12pt]{01_slides/img/tool.pdf}~The starting point\dots}
    \centering
    \includegraphics[width=\textwidth]{01_slides/img/jupyter.png}

    \vfill
    \begin{enumerate}
        \item Double click the \texttt{notebooks} directory
        \item Click on the \texttt{Python 3 (ipykernel)} notebook
    \end{enumerate}

    \begin{tikzpicture}[overlay]
        \node<2->[set red, align=center, text width=15em, xshift=5em] at (current page.west) {\textbf{CTRL + s} to save your work! {\footnotesize (Call the notebook \texttt{first\_tutorial.ipynb})}};
    \end{tikzpicture}
\end{frame}

\begin{frame}[fragile]{The Jupyter Interface}{Cells and execution\dots}
    The notebook is organised in \enquote{cells}:
    \begin{notebookin}
# Here goes your Python code
    \end{notebookin}

    \vfill

    You need to use \texttt{SHIFT + Enter} to execute the content of a cell:
    \begin{notebookin}
print("Hello, world!")
    \end{notebookin}
    \begin{notebookout}
Hello, world!
    \end{notebookout}
    \begin{notebookin}
3 + 2
    \end{notebookin}
    \begin{notebookout}
5
\end{notebookout}
\end{frame}

\begin{frame}[fragile]{The Jupyter Interface}{Cells and execution\dots}
    The top bar allows you to choose between \emph{Code} or \emph{Markdown}:
    \begin{center}
        \includegraphics[height=12pt]{01_slides/img/jupyter_bar.png}
        \hfil
        \includegraphics[height=12pt]{01_slides/img/jupyter_bar_md.png}
    \end{center}

    \vfill

    In Markdown mode, this cell
    \begin{notebookinmd}
# This is a title

This is text.

## This is a subtitle.

This is **bold** and this is *italic*.
    \end{notebookinmd}
    will display as text: useful for notes!
\end{frame}

\begin{frame}[fragile]{The Jupyter Interface}{Cells and execution\dots}
    You can use the commands on the right of the cell
    \begin{center}
        \includegraphics[height=12pt]{01_slides/img/jupyter_cell_cmd.png}
    \end{center}
    to
    \begin{itemize}
        \item \textbf{Duplicate} the current cell
        \item \textbf{Move} the current cell up or down
        \item \textbf{Add} a cell above or below the current cell
        \item \textbf{Delete} the current cell
    \end{itemize}
\end{frame}

\begin{frame}[fragile]{Some Basic Statistics}{\includegraphics[height=12pt]{01_slides/img/tool.pdf}~Back to the notebooks\dots}
    Start by activating some useful Jupyter extensions:
    \begin{notebookin}
%load_ext autoreload
%autoreload 2
    \end{notebookin}

    \vfill
    Then, load the libraries
    \begin{notebookin}
import numpy as np
import matplotlib.pyplot as plt
    \end{notebookin}

    \vfill
    We shall now take a look at some standard operations to perform on objects (\textbf{tensors}) in Python.
\end{frame}

\begin{frame}[fragile]{Some Basic Statistics}{\includegraphics[height=12pt]{01_slides/img/tool.pdf}~Some basic Python\dots}
    You can now perform some basic operations:
    \begin{notebookin}
# Generate 1000 samples in a standard normal distribution
X = np.random.randn(1000)
print("X.shape =", X.shape)
    \end{notebookin}
    \begin{notebookout}
X.shape = (1000,)
    \end{notebookout}

    \vfill \pause
    You can also manipulate \textbf{tensors} and \textbf{scalars}:
    \begin{notebookin}
# Generate 1000 samples in a normal distribution
# centred in 1.0 and with standard deviation 1/2
Y = 1.0 + 0.5 * np.random.randn(1000)
print("Y.shape =", Y.shape)
    \end{notebookin}
    \begin{notebookout}
Y.shape = (1000,)
    \end{notebookout}
\end{frame}

\begin{frame}[fragile]{Some Basic Statistics}{\includegraphics[height=12pt]{01_slides/img/tool.pdf}~Some basic Python\dots}
    You can reshape easily your tensors:
    \begin{notebookin}
X = X.reshape(-1, 1); Y = Y.reshape(-1, 1)
print("X.shape =", X.shape)
print("Y.shape =", Y.shape)
    \end{notebookin}
    \begin{notebookout}
X.shape = (1000, 1)
Y.shape = (1000, 1)
    \end{notebookout}

    \vfill \pause
    And perform basic operations (e.g. \textbf{matrix multiplication}):
    \begin{notebookin}
Z = (X.T @ Y) / 1000
print("Z.shape =", Z.shape)
print("Z =", Z.ravel())
    \end{notebookin}
    \begin{notebookout}
Z.shape = (1, 1)
Z = [0.0098006]
    \end{notebookout}
\end{frame}

\begin{frame}[fragile]{Some Basic Statistics}{\includegraphics[height=12pt]{01_slides/img/tool.pdf}~Some basic Python\dots}
    In general, you can manipulate \textbf{tensors} of any shape:
    \begin{notebookin}
# 1000 typical, 3-channel (RGB) images with square factor
X = np.random.randn(1000, 3, 64, 64)
print("X.shape =", X.shape)
    \end{notebookin}
    \begin{notebookout}
X.shape = (1000, 3, 64, 64)
    \end{notebookout}

    \vfill \pause
    And compute \textbf{statistics} on them:
    \begin{notebookin}
# The "axis" defines where to compute the operation
X_mean = X.mean(axis=1)  # in "torch", it will be "dim=1"
X_std = X.std(axis=1)
print("E[X].shape =", X_mean.shape)
print("std(X).shape = ", X_std.shape)
    \end{notebookin}
    \begin{notebookout}
E[X].shape = (1000, 64, 64)
std(X).shape =  (1000, 64, 64)
    \end{notebookout}
\end{frame}

\begin{frame}[fragile]{Some Basic Statistics}{\includegraphics[height=12pt]{01_slides/img/tool.pdf}~Some basic Python\dots}
    Tensors can be \textbf{sliced} and modified:
    \begin{notebookin}
# For instance:
#    X[3] takes the 4th element of a tensor
#    X[:, 2] takes all elements in the first axis, 3rd in the second
#    X[2:, :] takes from 3rd to the end in the first, all in the second
X = 3.0 * np.random.randn(128, 128)
X[32:64, 18:64] = 0
    \end{notebookin}

    \vfill \pause
    Certain objects (mostly 2d) can be shown:
    \begin{notebookin}
fig, ax = plt.subplots(ncols=1, nrows=1, layout="constrained")
img = ax.matshow(X, cmap="RdBu_r")
plt.colorbar(img)
plt.savefig("matplotlib_img.png")
    \end{notebookin}

    \begin{tikzpicture}[overlay]
        \node<3->[set, xshift=3em] at (current page.center) {\includegraphics[width=15em]{01_slides/img/matplotlib_img.png}};
    \end{tikzpicture}
\end{frame}

\begin{frame}[fragile]{Some Basic Statistics}{\includegraphics[height=12pt]{01_slides/img/tool.pdf}~Some basic Python\dots}
    We can also plot curves:
    \begin{notebookin}
# Create some data
x = np.linspace(0.0, 2*np.pi, num=1000)
y1 = np.sin(x)
y2 = np.cos(x)

# Create line plots
fig, ax = plt.subplots(ncols=1, nrows=1, layout="constrained")
ax.plot(x, y1, "r-", label="sin(x)")
ax.plot(x, y2, "b--", label="cos(x)")
ax.set(xlabel="x", ylabel="y")
ax.legend()
plt.savefig("matplotlib_lines.png")
    \end{notebookin}

    \begin{tikzpicture}[overlay]
        \node<2->[set, xshift=3em] at (current page.center) {\includegraphics[width=15em]{01_slides/img/matplotlib_lines.png}};
    \end{tikzpicture}
\end{frame}

\begin{frame}[fragile]{Some Basic Statistics}{\includegraphics[height=12pt]{01_slides/img/tool.pdf}~Some basic Python\dots}
    We can also plot histograms:
    \begin{notebookin}
# Create some data
x = np.random.beta(2.0, 5.0, size=2500)

# Create the histogram
fig, ax = plt.subplots(ncols=1, nrows=1, layout="constrained")
ax.hist(x, bins=100)
ax.set(xlabel="x", ylabel="Beta distribution")
plt.savefig("matplotlib_beta.png")
    \end{notebookin}

    \begin{tikzpicture}[overlay]
        \node<2->[set, xshift=3em, yshift=-5em] at (current page.center) {\includegraphics[width=15em]{01_slides/img/matplotlib_beta.png}};
    \end{tikzpicture}
\end{frame}

\begin{frame}[fragile]{Some Basic Statistics}{\includegraphics[height=12pt]{01_slides/img/tool.pdf}~Hands on, maybe\dots}
    Let us build a small 1D dataset:
    \begin{notebookin}
X_1 = np.random.randn(1500)
y_1 = np.zeros(1500)
X_2 = 3.0 + 0.25 * np.random.randn(1000)
y_2 = np.ones(1000)
X_3 = -4.1 + 2.0 * np.random.randn(500)
y_3 = 2 * np.ones(500)

X = np.concat([X_1, X_2, X_3], axis=0)
y = np.concat([y_1, y_2, y_3], axis=0).astype("int")

# Shuffle the entries
p = np.random.permutation(X.shape[0])
X = X[p]; y = y[p]
    \end{notebookin}

    \textbf{Goal:} Use $X$ as independent variables. Use $y$ as associated labels.
\end{frame}

\begin{frame}[fragile]{Some Basic Statistics}{\includegraphics[height=12pt]{01_slides/img/tool.pdf}~Hands on, maybe\dots}

    \begin{enumerate}
        \item \alert<2>{\includegraphics[height=12pt]{01_slides/img/pen-tool.pdf}~Verify that both $X$ and $y$ are vectors of length $3000$}
        \item \alert<3-4>{\includegraphics[height=12pt]{01_slides/img/pen-tool.pdf}~Plot the three univariate normal distributions (one per plot + on the same plot)}
        \item \alert<5>{\includegraphics[height=12pt]{01_slides/img/brain.pdf}~If $X$ are samples, what do labels identify in this example?}
    \end{enumerate}

    \vfill

    \begin{onlyenv}<2>
        \begin{notebookin}
print("X.shape =", X.shape)
print("y.shape =", y.shape)
        \end{notebookin}
        \begin{notebookout}
X.shape = (3000,)
y.shape = (3000,)
        \end{notebookout}
    \end{onlyenv}

    \begin{onlyenv}<3>
        \begin{tikzpicture}[overlay]
            \node[set, yshift=-12em] at (current page.center) {\includegraphics[width=0.85\textwidth]{01_slides/img/bayes_gauss.png}};
        \end{tikzpicture}
    \end{onlyenv}

    \begin{onlyenv}<4>
        \begin{notebookin}
fig, ax = plt.subplots(ncols=4, nrows=1, figsize=(20, 5), layout="constrained")
ax[0].hist(X_1, bins=50)
ax[0].set(xlabel="x", ylabel="y", title="Gaussian no. 1")
...  # same on ax[1] and ax[2]
ax[3].hist(X_1, bins=50, alpha=0.5, label="X_1")
...  # same but on same ax[3]
plt.savefig("bayes_gauss.png")
        \end{notebookin}
    \end{onlyenv}
\end{frame}

\begin{frame}[fragile]{Some Basic Statistics}{Hands on, maybe\dots}
    Now suppose we lost the original parameters of the three Gaussian.
    We are left with a \textbf{mixture}

    \begin{center}
        \includegraphics[width=0.3\textwidth]{01_slides/img/bayes_gauss_mix.png}
    \end{center}

    and with simple \textbf{labels} that just tell us which sample correspond to which Gaussian.

    \pause

    \begin{center}
        \includegraphics[height=12pt]{01_slides/img/brain.pdf}~Can we design an algorithm to determine the parameters and assign labels to new samples?
    \end{center}
\end{frame}

\begin{frame}[fragile]{Some Basic Statistics}{Hands on, maybe\dots}
    Recall an important lesson in statistics (\textbf{Bayes' theorem}):
    \begin{equation*}
        \underbrace{P(C \mid X)}_{\text{posterior}}
        \propto
        \underbrace{P(X \mid C)}_{\text{likelihood}}
        \underbrace{P(C)}_{\text{prior}}
    \end{equation*}

    \pause
    In other words,
    \begin{enumerate}[<+->]
        \item Given the probability of finding a certain \textbf{class} (label) $C = c$ {\tiny (favourable cases / total cases)}\dots
        \item \dots and the \enquote{model} (\textbf{likelihood}) which we use to describe our data (i.e.\ a Gaussian distribution) in a certain class,
        \item We can \enquote{update} the probability that $X = x$ is in a certain class by multiplying the two objects
    \end{enumerate}

    \pause
    The objective is thus:
    \begin{enumerate}[<+->]
        \item Determine (\enquote{fit}) correctly the parameters that describe the data (estimate mean and standard deviation for each class)
        \item \enquote{try} to apply each class to $X = x$ and see which class has the highest probability
    \end{enumerate}
\end{frame}

\begin{frame}[fragile]{Some Basic Statistics}{\includegraphics[height=12pt]{01_slides/img/tool.pdf}~Hands on, maybe\dots}
    \begin{enumerate}
        \item \alert<2>{\includegraphics[height=12pt]{01_slides/img/pen-tool.pdf}~Go in directory of your package in the \texttt{src} directory and create a new file named \texttt{naive\_classifier.py}}
        \item \alert<3>{\includegraphics[height=12pt]{01_slides/img/pen-tool.pdf}~Import \texttt{numpy} in the file and create the following functions:}
              \begin{enumerate}
                  \item \alert<4>{\texttt{gaussian\_distribution} to define a Gaussian probability of a sample $X$ given the mean and standard deviation of the curve}
                  \item \alert<5>{\texttt{fit} to separate the \enquote{training} samples into their classes and compute the mean and standard deviation of their distributions, and the \emph{prior}}
                  \item \alert<6>{\texttt{predict} to apply the Gaussian model for each possible class and select the one with the maximum \emph{posterior}}
                  \item \alert<7>{\texttt{accuracy} to compute how many times the prediction is exactly the same as the original label}
              \end{enumerate}
        \item \alert<8>{\includegraphics[height=12pt]{01_slides/img/pen-tool.pdf}~Go back to the notebook and import the \texttt{fit} and \texttt{predict} functions}
        \item \alert<9>{\includegraphics[height=12pt]{01_slides/img/pen-tool.pdf}~Apply the \texttt{fit} function to a subset of the $X$ dataset (e.g.\ the first 2500 elements)}
        \item \alert<10>{\includegraphics[height=12pt]{01_slides/img/pen-tool.pdf}~Apply the \texttt{predict} function to the remaining samples and compute the accuracy fo the predictions}
    \end{enumerate}

    \begin{onlyenv}<2>
        \begin{tikzpicture}[overlay]
            \node[label=right:src/] (src) at (current page.center) {\includegraphics[height=12pt]{01_slides/img/folder.pdf}};
            \node[below right=0.25em of src, label=right:mlintro/] (mod) {\includegraphics[height=12pt]{01_slides/img/folder.pdf}};
            \node[below right=0.25em of mod, label=right:\_\_init\_\_.py] (init) {\includegraphics[height=12pt]{01_slides/img/python.pdf}};
            \node[below=0.25em of init, label=right:naive\_classifier.py] (class) {\includegraphics[height=12pt]{01_slides/img/python.pdf}};

            \begin{scope}[on background layer]
                \node[set, fit=(src)(mod)(init)(class), minimum width=18em, minimum height=10em, xshift=7em, yshift=-3em] at (current page.center) {};
            \end{scope}
        \end{tikzpicture}
    \end{onlyenv}

    \begin{onlyenv}<3>
        \begin{tikzpicture}[overlay]
            \node[xshift=-2em, yshift=-5em] at (current page.center) {
                \begin{codewithcard}{naive_classifier.py}
import numpy as np
                \end{codewithcard}
            };
        \end{tikzpicture}
    \end{onlyenv}

    \begin{onlyenv}<4>
        \begin{tikzpicture}[overlay]
            \node[xshift=-2em, yshift=-5em] at (current page.center) {
                \begin{codewithcard}{naive_classifier.py}
...
def gaussian_distribution(x, mean, std):
    num = np.exp(-0.5 * (x-mean)**2 / std**2)
    den = np.sqrt(2.0 * np.pi) * std
    return num / den
                \end{codewithcard}
            };
        \end{tikzpicture}
    \end{onlyenv}

    \begin{onlyenv}<5>
        \begin{tikzpicture}[overlay]
            \node[xshift=-2em, yshift=-2em] at (current page.center) {
                \begin{codewithcard}{naive_classifier.py}
...
def fit(X, y):
    # Order the samples by their label
    idx = np.argsort(y)
    X_copy = X[idx]; y_copy = y[idx]
    # Recover the first positions of labels and split
    splits = np.unique(y_copy, return_index=True)[1]
    X_splits = np.split(X_copy, splits[1:])
    # Compute mean and std
    mean_list = []; std_list = []; prior_list = []
    for X_split in X_splits:
        mean_list.append(float(X_split.mean()))
        std_list.append(float(X_split.std()))
        prior_list.append(float(X_split.shape[0] / X.shape[0]))
    return mean_list, std_list, prior_list
                \end{codewithcard}
            };
        \end{tikzpicture}
    \end{onlyenv}

    \begin{onlyenv}<6>
        \begin{tikzpicture}[overlay]
            \node[xshift=-2em, yshift=-2em] at (current page.center) {
                \begin{codewithcard}{naive_classifier.py}
...
def predict(X, mean_list, std_list, prior_list):
    # Assume X is a list of values
    labels = []
    for x in X:
        # Compute the posterior for each class
        posterior_list = []
        for mean, std, prior in zip(mean_list, std_list, prior_list):
            likelihood = gaussian_distribution(x, mean, std)
            posterior = likelihood * prior
            posterior_list.append(posterior)
        labels.append(np.argmax(posterior_list))
    return np.array(labels).astype("int")
                \end{codewithcard}
            };
        \end{tikzpicture}
    \end{onlyenv}

    \begin{onlyenv}<7>
        \begin{tikzpicture}[overlay]
            \node[xshift=-2em, yshift=3em] at (current page.center) {
                \begin{codewithcard}{naive_classifier.py}
...
def accuracy(y_true, y_pred):
    return float((y_true == y_pred).mean())
                \end{codewithcard}
            };
        \end{tikzpicture}
    \end{onlyenv}

    \begin{onlyenv}<8>
        \begin{tikzpicture}[overlay]
            \node[set, xshift=-2em, yshift=3em] at (current page.center) {
                \begin{notebookin}
from mlintro.naive_classifier import fit, predict, accuracy
                \end{notebookin}
            };
        \end{tikzpicture}
    \end{onlyenv}

    \begin{onlyenv}<9>
        \begin{tikzpicture}[overlay]
            \node[set, xshift=-2em, yshift=3em] at (current page.center) {
                \begin{notebookin}
mean_list, std_list, prior_list = fit(X[:2500], y[:2500])
            \end{notebookin}
            };
        \end{tikzpicture}
    \end{onlyenv}

    \begin{onlyenv}<10>
        \begin{tikzpicture}[overlay]
            \node[set, text width=40em, xshift=-2em, yshift=3em] at (current page.center) {
                \begin{notebookin}
preds = predict(X[2500:], mean_list, std_list, prior_list)
acc = accuracy(y[2500:], preds)
print(f"Accuracy: {acc:.1%}")
                \end{notebookin}
                \begin{notebookout}
Accuracy: 96.6%
            \end{notebookout}
            };
        \end{tikzpicture}
    \end{onlyenv}
\end{frame}

\begin{frame}{Some Basic Statistics}{Bravo!}
    If you got to this point, then \emph{congratulations}! You just coded your first ML algorithm completely from scratch!

    \vspace{1em}

    What you just coded is called:
    \begin{center}
        \huge
        \textbf{Naive Bayes Classifier}
    \end{center}
    and performs simple classification, based on basic assumptions on the shape of the data (here: Gaussian classes). With several features, \emph{naive} means assuming the features are \textbf{independent} within each class.

    \vspace{1em}

    After breaking the ice, we can now proceed with the case study that will accompany us for the rest of the modules.
\end{frame}
